\documentclass[11pt]{article}

% ---------- JMLR-style layout (self-contained; no external jmlr2e.cls needed) ----------
\usepackage[letterpaper,margin=1in]{geometry}
\usepackage{times}
\usepackage{amsmath,amssymb}
\usepackage{booktabs}
\usepackage{array}
\usepackage{graphicx}
\usepackage{xcolor}
\usepackage{hyperref}
\usepackage{fancyhdr}
\usepackage{caption}
\usepackage{listings}
\usepackage{titlesec}

\hypersetup{colorlinks=true, linkcolor=blue, citecolor=blue, urlcolor=blue}

\titleformat{\section}{\normalfont\large\bfseries}{\thesection}{1em}{}
\titleformat{\subsection}{\normalfont\normalsize\bfseries}{\thesubsection}{1em}{}

\pagestyle{fancy}
\fancyhf{}
\fancyhead[L]{\footnotesize\slshape Timeout Enforcement and Reconciliation Report}
\fancyhead[R]{\footnotesize\slshape 2026}
\fancyfoot[C]{\footnotesize\thepage}
\renewcommand{\headrulewidth}{0.5pt}

\lstset{
  basicstyle=\ttfamily\footnotesize,
  breaklines=true,
  frame=single,
  columns=fullflexible,
  backgroundcolor=\color{gray!8},
  captionpos=b
}

% ---------- JMLR-style title block ----------
\newcommand{\jmlrtitle}[1]{
  \begin{center}
    {\LARGE\bfseries #1}
  \end{center}
  \vspace{0.3cm}
}

\begin{document}

\thispagestyle{empty}

\jmlrtitle{Timeout Enforcement, Configuration Reconciliation,\\ and Determinism Fixes in the Comparative Symbolic-Regression Benchmark Suite}

\begin{center}
\textbf{Investigation and Patch Report} \\[2pt]
\textit{Issues 2, 10a, 10b — Phase 3 Backlog} \\[6pt]
\today
\end{center}

\vspace{0.4cm}
\noindent\rule{\textwidth}{0.5pt}

\begin{abstract}
\noindent
This report documents the investigation and resolution of three open items
from the Phase 3 backlog of the comparative symbolic-regression benchmark
suite: (i) an unresolved configuration discrepancy between two write-ups
reporting 900\,s and 1100\,s PySR timeouts for the same sweep (Issue~10a);
(ii) a timeout-enforcement bug in which a stated 300\,s per-test limit
failed to prevent a 27{,}574\,s hang (Issue~10b); and (iii) a
per-process-randomized \texttt{hash()} seeding bug affecting five
unseeded or inconsistently-seeded neural-network code paths (Issue~2,
follow-up). We report root causes, the code-level fixes applied, and
benchmark results recomputed from the post-fix run logs across noiseless,
noisy, and noise-sweep protocols. Issue~10a is closed as a documentation
reconciliation; Issue~10b is closed with a process-group isolation fix to
\texttt{\_run\_pysr\_in\_subprocess}; Issue~2's follow-up patches five
call sites across two files, verified by static grep and successful
compilation.
\end{abstract}

\noindent\rule{\textwidth}{0.4pt}
\vspace{0.3cm}

\section{Introduction}

The Phase~3 backlog (Table~\ref{tab:backlog}) identified three
investigation-required items blocking closure of the benchmark
reproducibility audit. This report addresses all three, in order of
increasing severity: a stale documentation mismatch (10a), a silent
timeout-enforcement failure (10b), and a determinism bug affecting
reported pass rates (Issue~2 follow-up).

\begin{table}[h]
\centering
\caption{Phase 3 backlog items addressed in this report.}
\label{tab:backlog}
\begin{tabular}{@{}clc@{}}
\toprule
\textbf{\#} & \textbf{Item} & \textbf{Category} \\
\midrule
2   & 27--30/30 M4 noiseless pass-rate variance across 12 runs & Determinism (code) \\
10a & 900\,s vs.\ 1100\,s config bleed between two write-ups   & Documentation \\
10b & 300\,s timeout not enforced in harness code              & Enforcement bug (code) \\
\bottomrule
\end{tabular}
\end{table}

\section{Issue 10a: Timeout Configuration Reconciliation}
\label{sec:10a}

\subsection{Discrepancy}

\texttt{tab:sweeps} reported a PySR timeout of 1100\,s for the sweep run,
while the Phase~3 reproduction command block reported 900\,s for both
\texttt{method-timeout} and \texttt{pysr-timeout}. The two values were
never reconciled in the source document.

\subsection{Investigation}

Fourteen run-log JSON files were inspected for embedded invocation
metadata (Table~\ref{tab:runlogs}). None of the files retain the
configured timeout value, CLI string, or \texttt{argv} — only
\texttt{protocol.mode}, \texttt{noise\_level}, and \texttt{threshold} are
logged per run. As an indirect check, the maximum per-test wall-clock
time was extracted from every run; if either candidate value had actually
been the enforced ceiling, tests would cluster near it.

\begin{table}[h]
\centering
\caption{Run logs inspected for Issue 10a. No file contains a configured
timeout value; the closest available signal is per-test wall-clock time.}
\label{tab:runlogs}
\small
\begin{tabular}{@{}lccc@{}}
\toprule
\textbf{Run} & \textbf{Noise level} & \textbf{Threshold} & \textbf{Max test time (s)} \\
\midrule
\texttt{protocol\_core\_noiseless\_113233} & 0.0   & 0.999999 & 199.3 \\
\texttt{protocol\_core\_noisy\_115945}     & 0.1   & 0.90     & 251.0 \\
\texttt{protocol\_core\_noisy\_120017}     & 0.05  & 0.95     & 243.6 \\
\texttt{protocol\_core\_noisy\_120334}     & 0.01  & 0.99     & 249.0 \\
\texttt{protocol\_core\_noisy\_120417}     & 0.005 & 0.995    & 269.2 \\
\texttt{noise\_sweep\_*} (5 shards)        & 0.0--0.1 & mixed & 140.1 (mean, slowest method) \\
\bottomrule
\end{tabular}
\end{table}

The highest observed per-test time across all 5 protocol-core runs and 5
sweep shards is 269.2\,s — well short of \emph{both} 900\,s and 1100\,s.
Every run in the dataset shows \texttt{n\_success == n\_total}: zero
timeouts occurred regardless of which value was actually configured. The
dataset is therefore \textbf{silent} on which of the two numbers is
correct; both were sufficient headroom for the observed run times.

\subsection{Resolution}

Because the two values cannot be adjudicated from any available log, and
because neither value changes any reported result, the fix standardizes
on the one number that is independently verifiable — the \texttt{pysr-timeout
900} in the Phase~3 reproduction command, which is a literal, copyable,
runnable command. The 1100\,s figure in \texttt{tab:sweeps} is treated as
a stale draft artifact and removed.

\begin{quote}
\textbf{Before} (\texttt{tab:sweeps} note): NN seeds~=~5, random seed~42,
method timeout~=~900\,s, \textbf{PySR timeout~=~1100\,s}, threshold
(noisy)~=~0.995, threshold (noiseless)~=~0.999999.

\textbf{After}: NN seeds~=~5, random seed~42, method timeout~=~900\,s,
\textbf{PySR timeout~=~900\,s}, threshold (noisy)~=~0.995, threshold
(noiseless)~=~0.999999.
\end{quote}

\textbf{Status: closed (documentation only).} If original job-scheduler
logs surface later showing 1100\,s was in fact used, this should be
reopened; the current fix is a best-available reconciliation rather than a
confirmed root cause.

\section{Issue 10b: Timeout Enforcement Failure}
\label{sec:10b}

\subsection{Symptom}

The reproducibility statement records a 300\,s \texttt{method-timeout} for
tests 1--18, yet Arrhenius (test 4) ran for 27{,}574\,s — over $90\times$
the configured limit — before the underlying Julia process was killed
manually. Nernst (test 6) also timed out under the same 300\,s default.

\subsection{Root cause}

Two independent bugs were found in \texttt{\_run\_pysr\_in\_subprocess}
and its surrounding timeout scaffolding:

\begin{description}
\item[Bug 1 (root cause).] The PySR worker is spawned with
\texttt{subprocess.Popen(\ldots)} without process-group isolation
(\texttt{start\_new\_session} not set). On \texttt{TimeoutExpired},
\texttt{proc.kill()} sends \texttt{SIGKILL} only to the direct Python
wrapper child. If that wrapper in turn launches Julia as its own OS
subprocess, killing the wrapper orphans Julia, which has no timeout of
its own and keeps running until manually noticed — consistent with the
documented wording ``killed'' rather than ``timed out.''
\item[Bug 2 (dead safety net).] A hard-kill fallback calls
\texttt{\_kill\_thread(...)}, referencing a function that is never
defined anywhere in the file. Every invocation raises
\texttt{NameError}, silently swallowed by a blanket
\texttt{except~Exception:~pass}, making the intended last-resort kill a
permanent no-op.
\end{description}

\subsection{Fix applied}

\begin{lstlisting}[language=Python,caption={Process-group isolation fix applied to \texttt{\_run\_pysr\_in\_subprocess}.}]
proc = subprocess.Popen(
    [sys.executable, "-c", _SUBPROCESS_WORKER],
    stdin=subprocess.PIPE, stdout=subprocess.PIPE, stderr=subprocess.PIPE,
    env=env,
    start_new_session=True,   # new process group
)
...
except subprocess.TimeoutExpired:
    _kill_process_group(proc)   # SIGKILL whole group; psutil fallback
    proc.wait()
\end{lstlisting}

Bug 2's dead \texttt{\_kill\_thread} call was removed rather than
implemented, since the real hang is an orphaned OS process, not a blocked
Python thread — killing the thread would not have touched Julia either.

\subsection{Verification plan}

\begin{enumerate}
\item Re-run Arrhenius (test 4) alone with a short \texttt{--method-timeout}.
\item Inspect \texttt{ps -ef --forest} (or \texttt{psutil} children) for an
orphaned \texttt{julia} process after the wrapper is killed.
\item Confirm the whole process group terminates within the configured
window.
\item Re-run tests 1--18 with the fixed harness and flag that earlier
results from that window may have suffered silent resource contention
(e.g.\ Nernst/test 6 competing with an orphaned Arrhenius process for CPU).
\end{enumerate}

\textbf{Status: code fix applied to
\texttt{run\_comparative\_suite\_benchmark\_v2\_FIXED.py}; end-to-end
re-run verification pending.}

\section{Issue 2 Follow-Up: Unseeded Neural-Network Determinism}
\label{sec:issue2}

\subsection{Root cause}

Python's built-in \texttt{hash()} on a \texttt{str} is randomized
per-process unless \texttt{PYTHONHASHSEED} is pinned. Code seeding an RNG
with \texttt{hash(description)} under the assumption of reproducibility
was in fact producing a different seed on every run or shard. Separately,
\texttt{ImprovedNNMethod.run()} — the library's own default entry point —
performed no seeding at all.

\subsection{Fixes applied}

Five call sites across two files were patched, replacing
\texttt{hash()}/no seeding with \texttt{hashlib.sha256}-derived seeds
(Table~\ref{tab:patchsites}).

\begin{table}[h]
\centering
\caption{Five call sites patched for Issue 2 follow-up.}
\label{tab:patchsites}
\small
\begin{tabular}{@{}clc@{}}
\toprule
\textbf{\#} & \textbf{Location} & \textbf{Before} \\
\midrule
1 & \texttt{BaseMethod.\_nn\_residual\_fit} & unseeded \\
2 & \texttt{HybridAllDomainsMethod.\_nn\_residual\_fit} (own override) & unseeded \\
3 & \texttt{HybridDeFiMethod} reconstruction-fallback seed & \texttt{hash()} \\
4 & Noise-injection seed in \texttt{main()} & \texttt{hash()} \\
5 & \texttt{ImprovedNNMethod.run()} default (\texttt{nn\_seeds=1}) path & unseeded \\
\bottomrule
\end{tabular}
\end{table}

Fix 5 required a guard (\texttt{if self.\_nn\_seeds == 1}) so the new
seeding does not disturb \texttt{run\_multiseed()}'s independent
per-trial seeding when \texttt{nn\_seeds > 1}.

\subsection{Verification}

\texttt{python3 -m py\_compile} passed with no syntax errors; a
non-comment grep for \texttt{hash(} across the patched file returned zero
matches; both \texttt{\_nn\_residual\_fit} definitions and call sites were
confirmed consistent on the new \texttt{description} parameter.

\textbf{Impact.} This is a behavior-changing fix. All prior runs of this
script — including results behind the five-system comparison tables —
were generated with these five paths unseeded or inconsistently seeded.
Re-running is expected to produce \emph{different} numbers, not merely
reproducible copies of the old ones.

\section{Benchmark Results (Post-Fix Run Logs)}
\label{sec:results}

The following tables are recomputed directly from the fourteen uploaded
post-fix run logs, applying each run's own logged threshold.

\subsection{Noiseless and noisy protocol-core runs}

\begin{table}[h]
\centering
\caption{Pass rate (equations solved above threshold, out of 30) by
method and noise condition. Threshold is set per the noise level
following the harness's own noise-appropriate calibration.}
\label{tab:core}
\small
\begin{tabular}{@{}lcccccc@{}}
\toprule
\textbf{Method} & \textbf{Noiseless} & \textbf{$\sigma$=0.005} & \textbf{$\sigma$=0.01} & \textbf{$\sigma$=0.05} & \textbf{$\sigma$=0.1} \\
 & (thr.\ 0.999999) & (thr.\ 0.995) & (thr.\ 0.99) & (thr.\ 0.95) & (thr.\ 0.9) \\
\midrule
PureLLM Baseline           & 29/30 & 30/30 & 29/30 & 30/30 & 30/30 \\
ImprovedNN                 &  0/30 & 27/30 & 28/30 & 29/30 & 29/30 \\
EnhancedHybridSystemDeFi   & 23/30 & 30/30 & 30/30 & 30/30 & 30/30 \\
HybridSystemLLMNN (all-domains) & \textbf{30/30} & \textbf{30/30} & \textbf{30/30} & \textbf{30/30} & \textbf{30/30} \\
SymbolicEngineWithLLM      & 28/30 & 29/30 & 29/30 & 30/30 & 30/30 \\
HybridDiscoverySystem v50\_2 & 21/30 & 23/30 & 26/30 & 27/30 & 25/30 \\
\bottomrule
\end{tabular}
\end{table}

\begin{table}[h]
\centering
\caption{Maximum per-test wall-clock time (seconds) observed in each run,
by method. Used as the indirect timeout check in Section~\ref{sec:10a};
no value approaches either 900\,s or 1100\,s.}
\label{tab:times}
\small
\begin{tabular}{@{}lccccc@{}}
\toprule
\textbf{Method} & \textbf{Noiseless} & \textbf{$\sigma$=0.005} & \textbf{$\sigma$=0.01} & \textbf{$\sigma$=0.05} & \textbf{$\sigma$=0.1} \\
\midrule
PureLLM Baseline           & 11.1 & 12.9 & 14.0 & 11.8 & 10.2 \\
ImprovedNN                 &  4.1 &  4.1 &  4.0 &  4.2 &  4.1 \\
EnhancedHybridSystemDeFi   & 13.1 & 14.0 & 13.9 & 15.2 & 14.5 \\
HybridSystemLLMNN (all-domains) & 14.4 & 12.6 & 13.0 & 15.9 & 12.6 \\
SymbolicEngineWithLLM      & 61.1 & 54.9 & 55.0 & 54.0 & 60.1 \\
HybridDiscoverySystem v50\_2 & 199.3 & \textbf{269.2} & 249.0 & 243.6 & 251.0 \\
\bottomrule
\end{tabular}
\end{table}

\subsection{Noise-sweep summary (\texttt{tab:sweeps}, post-reconciliation)}

\begin{table}[h]
\centering
\caption{Noise-sweep recovery rate and median $R^2$ by method and noise
level (5 shards, 30 equations each, \texttt{NN seeds = 5}). PySR timeout
reconciled to 900\,s per Section~\ref{sec:10a}.}
\label{tab:sweeps}
\small
\begin{tabular}{@{}lccccc@{}}
\toprule
& \multicolumn{5}{c}{\textbf{Noise level}} \\
\cmidrule(lr){2-6}
\textbf{Method} & \textbf{0.0} & \textbf{0.005} & \textbf{0.01} & \textbf{0.05} & \textbf{0.1} \\
\midrule
EnhancedHybridSystemDeFi   & 0.767 & 1.000 & 1.000 & 1.000 & 1.000 \\
HybridDiscoverySystem v50\_2 & 0.700 & 0.767 & 0.867 & 0.900 & 0.833 \\
HybridSystemLLMNN (all-domains) & \textbf{1.000} & \textbf{1.000} & \textbf{1.000} & \textbf{1.000} & \textbf{1.000} \\
ImprovedNN                 & 0.000 & 0.900 & 0.933 & 0.967 & 0.967 \\
PureLLM Baseline           & 0.967 & 1.000 & 0.967 & 1.000 & 1.000 \\
SymbolicEngineWithLLM      & 0.933 & 0.967 & 0.967 & 1.000 & 1.000 \\
\bottomrule
\end{tabular}

\vspace{4pt}
{\small All 5 shards show \texttt{n\_success == n\_total} for every
method (0 timeouts, 0 catastrophic failures reaching \texttt{n\_total}),
consistent with the ``all 30 equations completed with zero timeouts''
claim in \texttt{tab:sweeps} independent of the 900\,s/1100\,s question.}
\end{table}

\subsection{Pre-fix baseline (Feynman, random 80/20 split)}

\begin{table}[h]
\centering
\caption{Pre-fix baseline solve rate reported consistently across 4
rescued shards, compared against the paper's stated claim.}
\label{tab:baseline}
\begin{tabular}{@{}lcc@{}}
\toprule
\textbf{Quantity} & \textbf{Rescued shards (4/4 agree)} & \textbf{Paper claim} \\
\midrule
Pass count / total & 74/90 & --- \\
Solve rate & 0.822 & 9/30 = 0.300 \\
\bottomrule
\end{tabular}

\vspace{4pt}
{\small The rescued shards report a random 80/20-split protocol
(\texttt{n\_total = 90}), not directly comparable in denominator to the
paper's stated 9/30 (30-equation, non-split protocol); this discrepancy
is protocol-mismatch, not a numerical error, and is out of scope for this
report.}
\end{table}

\section{Summary of Status}

\begin{table}[h]
\centering
\caption{Status of all items addressed in this report.}
\begin{tabular}{@{}clc@{}}
\toprule
\textbf{\#} & \textbf{Item} & \textbf{Status} \\
\midrule
10a & 900\,s vs.\ 1100\,s config bleed & \textbf{Closed} (doc reconciliation) \\
10b & 300\,s timeout not enforced & \textbf{Code fix applied}; re-run verification pending \\
2 (follow-up) & Unseeded NN determinism (5 call sites) & \textbf{Patched \& verified} (compile + grep) \\
\bottomrule
\end{tabular}
\end{table}

\noindent
\textbf{Recommended next actions:} (1) re-run tests 1--18 with the
process-group fix and confirm no orphaned Julia processes remain after a
short forced timeout; (2) regenerate the five-system comparison tables
using the Issue-2-patched harness, keeping pre-patch outputs for a
before/after diff rather than discarding them; (3) if original
job-scheduler logs for the sweep run are located, revisit the 900\,s vs.\
1100\,s reconciliation in Section~\ref{sec:10a} against that ground
truth.

\end{document}
